\section{A quantitative inverse for the small-cut estimate}\label{sec:inverse}

This section refines the small-cut argument from the source manuscript.
It gives both a pair-structure conclusion and a concentration conclusion.
It does not give a decomposition of every graph with many cycles.

Let $G$ be a connected multigraph.
Let $A_1,\ldots,A_k$ be disjoint cyclic regions, where $k\ge2$ and $|\delta A_i|\le d$.
Set $\Delta=\max\{d-2,0\}$ and $K=1+\Delta$.
For a uniform even subgraph $\xi$, define
\[
 E_i=\{\xi\cap\delta A_i=\varnothing,\quad
             \xi\cap E(G[A_i])\ne\varnothing\},\qquad
 p_i=\Pr(E_i),\quad Z=\sum_i\mathbf1_{E_i},\quad \mu=\sum_i p_i.
\]
Let $\ell_{ij}$ count the links between $A_i$ and $A_j$.
The source manuscript gives
\[
 \Pr(E_i\cap E_j)=2^{\ell_{ij}-1}p_i p_j.
\]
Every region is proper, since $k\ge2$.
The cut formula therefore improves the elementary lower bound to
\begin{equation}\label{eq:pi-lower}
 p_i=2^{\operatorname{comp}(G-A_i)-|\delta A_i|}
       (1-2^{-\beta(G[A_i])})\ge2^{-d},\qquad \mu\ge k2^{-d}.
\end{equation}

\begin{theorem}\label{thm:inverse}
Suppose that $0\le\eta<1/2$ and
\[
 f:=\Pr\bigl(\xi\cap E(G[A_i])\text{ is a forest for every }i\bigr)
       \ge\tfrac12-\eta.
\]
Define the ordered one-link mass by
\[
 D_1=\frac1{\mu^2}\sum_{i\ne j:\,\ell_{ij}=1}p_i p_j.
\]
Then
\begin{equation}\label{eq:inverse-mass}
 D_1\le \frac{4\eta}{1+2\eta}+\frac K\mu
       \le \frac{4\eta}{1+2\eta}+\frac{K2^d}{k}.
\end{equation}
Choose $I,J$ independently with $\Pr(I=i)=\Pr(J=i)=p_i/\mu$.
Then
\begin{equation}\label{eq:two-link-mass}
 \Pr(I\ne J,\ \ell_{IJ}=2)
   \ge1-\frac{4\eta}{1+2\eta}-\frac{2K}{\mu}.
\end{equation}
Writing $\nu=\Pr(Z>0)$, we also have
\begin{gather}
 \frac1{2+K/\mu}\le\nu\le\tfrac12+\eta,\label{eq:nu}\\
 \frac{\operatorname{Var}(Z\mid Z>0)}{(\mu/\nu)^2}
       \le2\eta+(\tfrac12+\eta)\frac K\mu.\label{eq:occupancy}
\end{gather}
\end{theorem}

\begin{proof}
The exceptional-partner lemma bounds each region's partners with at least three links by $\Delta$.
For these pairs, use $\Pr(E_i\cap E_j)\le p_i$.
For one-link pairs, the product moment is $p_i p_j$.
For two-link pairs, it is $2p_i p_j$.
Together with the diagonal terms, this gives
\begin{equation}\label{eq:inverse-moment}
 \mathbb E Z^2\le2\mu^2-D_1\mu^2+K\mu.
\end{equation}
On the event defining $f$, every $E_i$ fails.
Thus $\Pr(Z=0)\ge f$, and $\nu\le1/2+\eta$.
Cauchy--Schwarz gives $\mathbb E Z^2\ge\mu^2/\nu$.
Substitution in \eqref{eq:inverse-moment} proves \eqref{eq:inverse-mass}.

The weighted diagonal mass is at most $1/\mu$, since $p_i\le1$.
The weighted mass of exceptional ordered pairs is at most $\Delta/\mu$.
Connectedness gives $\ell_{ij}\ge1$ for distinct regions.
Subtract these masses and \eqref{eq:inverse-mass} from one to obtain \eqref{eq:two-link-mass}.

Dropping $D_1$ from \eqref{eq:inverse-moment} gives the lower bound in \eqref{eq:nu}.
Finally, $\mathbb E(Z\mid Z>0)=\mu/\nu$, and
\[
 \frac{\operatorname{Var}(Z\mid Z>0)}{(\mu/\nu)^2}
   =\frac{\nu\mathbb E Z^2}{\mu^2}-1
   \le2\nu-1+\nu K/\mu.
\]
The upper bound on $\nu$ proves \eqref{eq:occupancy}.
\end{proof}

For fixed $d$, let $k\to\infty$ and $\eta\to0$.
The theorem forces almost all weighted pair mass onto distinct two-link pairs.
It also forces $Z/\mu$ to concentrate near two, conditionally on being positive.
That conditioning event has probability tending to one half.
These conclusions describe the near-equality mechanism of this particular second-moment argument.

\paragraph{Transfer from a cycle count.}
The chord construction gives $c(G)/2^r\le p(G^+,W)+2^{-t}$ when $t$ edges are selected.
Hence $c(G)\ge(1/2-\eta)2^r$ implies $p(G^+,W)\ge1/2-\eta-2^{-t}$.
Suppose that the regions above are disjoint from $V(W)$ in $G^+$.
The forest event outside $W$ implies the event defining $f$.
The theorem therefore applies with $\eta+2^{-t}$, when that number is less than one half.
A separate argument must produce many regions with bounded cuts.

\paragraph{A more general second-moment boundary.}
If a proposed improvement uses only $\mathbb E Z=\mu$ and $\mathbb E Z^2\le2\mu^2+K\mu$, it cannot force an error exponentially small in $\mu$.
For $K>0$, take $Z=0$ with probability $(\mu+K)/(2\mu+K)$ and $Z=2\mu+K$ otherwise.
Both moment conditions hold with equality, and
\[
 \Pr(Z=0)-\tfrac12=\frac{K}{2(2\mu+K)}.
\]
This is a limitation of those moment inequalities alone.
It is not a graph construction or an obstruction to stronger graph estimates.
